\section{Sparse autoencoders}
\label{sec:sae}

\Cref{sec:features} ended with a problem: given many activation vectors, find a large set of directions such
that every vector is a combination of a few of them. The \emph{sparse autoencoder} (SAE) is the tool that the
field currently uses for this. It was introduced to interpretability by \citet{bricken2023monosemanticity} and
\citet{cunningham2024sparse} and has since been scaled to production models
\citep{templeton2024scaling,gao2025scaling}.

\subsection{Definition}

We fix a place in the model, for instance the residual stream after layer 6, and collect its activation
$x\in\R^d$ at every token of a large text corpus. An SAE with $m$ \emph{latents}, where $m$ is much larger than
$d$, consists of an encoder and a decoder,
\begin{equation}\label{eq:sae}
f(x) \;=\; \relu\big(\Wenc\,(x-\bdec)+\benc\big)\in\R^m,\qquad
\hat x(x) \;=\; \Wdec\,f(x)+\bdec\in\R^d ,
\end{equation}
with $\Wenc\in\R^{m\times d}$, $\Wdec\in\R^{d\times m}$, $\benc\in\R^m$ and $\bdec\in\R^d$. We write $\Wdec[j]\in\R^d$
for column $j$ of the decoder and call it the \emph{dictionary direction} of latent $j$. The columns are kept at
unit norm. The parameters are trained to minimise
\begin{equation}\label{eq:saeloss}
\text{loss} \;=\; \E\Big[\big\|x-\hat x(x)\big\|^2+\lambda\sum_{j=1}^{m}f_j(x)\Big],\qquad\lambda>0,
\end{equation}
where the expectation is over the collected activations. The first term asks for a faithful reconstruction.
The second term is the sum of the absolute values of the entries of $f(x)$, since $f_j(x)\ge0$. It is known as
an L1 penalty, and it pushes most latents to zero.
\Cref{fig:sae} (left) shows the architecture.

Comparing $\hat x(x)=\bdec+\sum_jf_j(x)\,\Wdec[j]$ with \eqref{eq:lrh}, we see what is hoped for: the dictionary
directions $\Wdec[j]$ should be the feature directions $v_j$, and the latent activations should be the feature
activations. We keep the two notions apart by name. A \emph{latent} is a unit of an SAE, which we can compute.
A \emph{feature} is a property that the model represents, which we want to know. Whether a given latent
corresponds to a feature is an empirical question in every single case.

The method descends from sparse coding in neuroscience \citep{olshausen1997sparse}. Classical sparse coding
solves a separate optimisation problem for the coefficients of each data point. An SAE replaces this by a single
learned encoder, which makes it cheap enough for billions of activations. Dictionary sizes in published work
range from a few thousand latents to $2^{25}\approx34$ million for a production model
\citep{templeton2024scaling}.

\begin{figure}[t]
\centering
\begin{tikzpicture}[>=Latex, font=\small]
\foreach \y/\c in {0/60,1/25,2/80,3/40,4/15} {\draw[tutgrey, fill=tutblue!\c] (0,0.34*\y) rectangle +(0.34,0.34);}
\node[below, align=center] at (0.17,-0.05) {$x\in\R^d$};
\foreach \y in {0,...,12} {\draw[tutgrey, fill=white] (3.1,0.34*\y-1.36) rectangle +(0.34,0.34);}
\draw[tutgrey, fill=tutred!80] (3.1,0.34*3-1.36) rectangle +(0.34,0.34);
\draw[tutgrey, fill=tutred!35] (3.1,0.34*9-1.36) rectangle +(0.34,0.34);
\node[below, align=center] at (3.27,-1.41) {$f(x)\in\R^m$\\ few nonzero};
\foreach \y/\c in {0/55,1/30,2/75,3/40,4/20} {\draw[tutgrey, fill=tutblue!\c] (6.2,0.34*\y) rectangle +(0.34,0.34);}
\node[below, align=center] at (6.37,-0.05) {$\hat x(x)\in\R^d$};
\draw[->, tutgrey] (0.5,0.85) -- node[above] {encoder} (2.95,0.85);
\draw[->, tutgrey] (3.6,0.85) -- node[above] {decoder} (6.05,0.85);
\end{tikzpicture}
\hfill
\begin{tikzpicture}
\begin{axis}[width=6.0cm, height=4.7cm, xmin=0, xmax=1, ymin=0, ymax=1, axis lines=left,
  xlabel={true activation $f^\star$}, ylabel={optimal latent activation}, label style={font=\small},
  tick label style={font=\footnotesize}, xtick={0,0.5,1}, ytick={0,0.5,1}, clip=false]
\addplot[dashed, tutgrey, domain=0:1, samples=2] {x};
\addplot[tutred, very thick, domain=0:1, samples=201] {max(0,x-0.2)};
\node[font=\footnotesize, tutgrey, anchor=south east] at (axis cs:0.74,0.76) {no shrinkage};
\node[font=\footnotesize, tutred!80!black, anchor=north west] at (axis cs:0.62,0.40) {$\max(0,f^\star-\lambda/2)$};
\end{axis}
\end{tikzpicture}
\caption{Left: a sparse autoencoder maps an activation to a much longer vector with few nonzero entries and
back. Right: the latent activation that minimises reconstruction error plus L1 penalty
(\cref{prop:shrink}, drawn for $\lambda=0.4$) lies below the true activation by $\lambda/2$.}
\label{fig:sae}
\end{figure}

\subsection{What the sparsity penalty does to the activations}

The penalty in \eqref{eq:saeloss} is a compromise. What we want to keep small is the \emph{number} of active
latents. That number is not differentiable, and the sum of the activations is used in its place. The price of
the substitute can be computed exactly in the simplest situation.

\begin{proposition}[Shrinkage]\label{prop:shrink}
Let $v\in\R^d$ be a unit vector and let $x=f^\star v$ with $f^\star\ge0$. Among all $f\ge0$, the value that
minimises $\|x-fv\|^2+\lambda f$ is
\[
f \;=\; \max\big(0,\ f^\star-\lambda/2\big).
\]
\end{proposition}
\begin{proof}
Since $\|v\|=1$, we have $\|x-fv\|^2=\|(f^\star-f)\,v\|^2=(f^\star-f)^2$. We therefore minimise
$(f^\star-f)^2+\lambda f$ over $f\ge0$. This function of $f$ is convex, and its derivative is
\begin{equation}\label{eq:shrinkderiv}
\frac{d}{df}\Big((f^\star-f)^2+\lambda f\Big) \;=\; -2\,(f^\star-f)+\lambda .
\end{equation}
If $f^\star\ge\lambda/2$, the derivative \eqref{eq:shrinkderiv} vanishes at $f=f^\star-\lambda/2\ge0$, which is
then the minimiser. If $f^\star<\lambda/2$, the derivative \eqref{eq:shrinkderiv} is at least
$\lambda-2f^\star>0$ for all $f\ge0$, so the function is increasing on $[0,\infty)$ and the minimiser is $f=0$.
\end{proof}

So a latent that fires for the right reason is still reported too small, by $\lambda/2$, and weak features are
silenced completely (\cref{fig:sae}, right). The reconstruction $\hat x$ is systematically too short. This effect
is called shrinkage or feature suppression \citep{wright2024addressing}. Current architectures avoid it by
separating the decision \emph{whether} a latent is active from the estimate of \emph{how strongly}.

\begin{itemize}
\item \textbf{TopK} \citep{gao2025scaling}. The $\relu$ in \eqref{eq:sae} is replaced by an operation that keeps
the $k$ largest entries and sets the others to zero. The number of active latents is then exactly $k$, and the
penalty is dropped.
\item \textbf{JumpReLU} \citep{rajamanoharan2024jumping}. Each latent has a learned threshold $\theta_j>0$.
An entry passes unchanged if it exceeds $\theta_j$ and is set to zero otherwise. The loss penalises the number
of active latents directly, and the missing gradient is replaced by an estimate.
\item \textbf{Gated} \citep{rajamanoharan2024improving}. One part of the encoder decides which latents are
active, and another part estimates their magnitudes.
\end{itemize}

\subsection{How to judge a dictionary}
\label{sec:saeeval}

An SAE has no ground truth to be compared with, since the true features are unknown. It is judged by several
numbers at once.

\begin{enumerate}
\item \textbf{Reconstruction.} The fraction of variance unexplained,
$\E\|x-\hat x(x)\|^2\,/\,\E\|x-\E x\|^2$, should be small.
\item \textbf{Sparsity.} The average number of nonzero latents per activation, usually called the L0 of the SAE, should be
small compared with $d$. Published dictionaries have values from about ten to a few hundred.
\item \textbf{Effect on the model.} One runs the model with $x$ replaced by $\hat x(x)$ and measures how much the
next-token loss increases. This matters because the model does not care equally about all directions of $x$. A
small reconstruction error can sit in exactly the direction that a later layer reads.
\item \textbf{Health.} \emph{Dead} latents never fire, and \emph{dense} latents fire on a large share of all
tokens. Both are wasted capacity.
\item \textbf{Interpretability.} For a sample of latents, a human or a language model writes a short
description of the texts on which the latent fires. The description is scored by how well it predicts the
activations on new texts \citep{bills2023language}.
\end{enumerate}

Items 1 and 2 trade off against each other, so two SAEs can only be compared along a whole curve.
\citet{karvonen2025saebench} collect these and further tests in a benchmark.

\subsection{Reading a single latent}
\label{sec:latent}

Suppose a latent has passed these tests. Three views are combined to understand it.

\begin{itemize}
\item \textbf{Where it fires.} One reads text on which the latent is active. The examples should be sampled over
the whole range of activations. The strongest ones alone give too clean a picture.
\item \textbf{What it writes.} By \cref{prop:dla}, adding a multiple of $\Wdec[j]$ to the final stream changes
the logit of token $t$ in proportion to $W_U[t]^\top(\Wdec[j]-\mean(\Wdec[j])\,\one)$. Sorting the vocabulary by
this number shows which tokens the latent promotes directly.
\item \textbf{What happens when it is changed.} One runs the model while holding $f_j$ at zero or at a large
value and adds the corresponding change of $\hat x$ to the stream. This is an intervention, and it is the only
one of the three views that tests whether the model uses the direction.
\end{itemize}

The best-known example is from \citet{templeton2024scaling}, who trained SAEs with up to $2^{25}$ latents on a
middle layer of the production model Claude 3 Sonnet. One latent fired on mentions and pictures of the Golden
Gate Bridge. When the authors clamped it to ten times its maximal observed value, the model began to describe
itself as the Golden Gate Bridge. Open dictionaries exist for open models, for example Gemma Scope for all layers
of Gemma 2 \citep{lieberum2024gemma}, and can be browsed on Neuronpedia \citep{lin2023neuronpedia}.

\subsection{What a sparse autoencoder does not give}

The appeal of SAEs is that they produce units with readable descriptions without any labels. The following
limits are established by now and have to be kept in mind when using them.

\begin{caution}
\textbf{The reconstruction is imperfect.} An analysis in terms of latents describes $\hat x$, and the remainder
$x-\hat x$ often matters for the output (item 3 of \cref{sec:saeeval}).

\textbf{The dictionary is not canonical.} The latents depend on the dictionary size. A latent of a small SAE
tends to split into several more specific latents in a larger one \citep{bricken2023monosemanticity}, and no
size is the correct one \citep{leask2025sparse}. They also depend on the random seed
\citep{paulo2025sparse}. A general latent can fail to fire on cases that a more specific latent has taken over,
which is called absorption \citep{chanin2024absorption}.

\textbf{Readable latents do not prove learned structure.} \citet{heap2025sparse} trained SAEs on transformers
with random weights and obtained latents with similar automatic interpretability scores. Part of the
sparse structure comes from the input text and not from anything the network has learned.

\textbf{Usefulness has to be shown per task.} On supervised tasks with labels, simple baselines are hard to
beat. \citet{kantamneni2025sparse} found no consistent advantage of probes on SAE latents over probes on raw
activations, and \citet{wu2025axbench} found that steering with SAE latents was weaker than simple alternatives.
\end{caution}

None of this makes SAEs useless. They remain the main tool for discovering concepts that one did not think of in
advance. It does mean that a latent is a hypothesis about the model, to be tested with the interventions of
\cref{sec:interventions}.

\subsection{Relatives: transcoders and crosscoders}
\label{sec:transcoders}

An SAE reconstructs an activation from itself. Two variants change what is reconstructed from what.
A \emph{transcoder} \citep{dunefsky2024transcoders} has the same form as \eqref{eq:sae}, and it is trained to
predict the \emph{output} of an MLP layer from the \emph{input} of that layer. It is a sparse, wide replacement
of the MLP. Its use is in circuit analysis. Between an SAE latent in one layer and an SAE latent in a later
layer lies a nonlinear MLP, and their interaction is hard to describe. With transcoders, the path from one
latent to the next passes only through linear maps once the attention patterns and the scales in $\LN$ are
held fixed, and it can be read from the weights. \Cref{sec:featurecircuits} builds on this. A \emph{crosscoder} \citep{lindsey2024crosscoders}
uses one set of latents to reconstruct several activations at once, either several layers of one model or the
same layer of two models. The second use allows one to ask what changed between a model and its fine-tuned
version.

\begin{takeaways}
A sparse autoencoder rewrites an activation as a sparse, non-negative combination of many learned directions.
The L1 penalty biases activations downward, which TopK and JumpReLU variants avoid. A dictionary is judged
by reconstruction, sparsity, effect on the model and interpretability together. Its latents are candidate
features. They depend on training choices, they leave an unexplained remainder, and they need causal tests.
\end{takeaways}

\begin{exercise}\label{ex:normconstraint}
Consider the loss \eqref{eq:saeloss} without the constraint that the columns of $\Wdec$ have unit norm. Show
that multiplying column $j$ of $\Wdec$ by $2$ and dividing row $j$ of $\Wenc$ and entry $j$ of $\benc$ by $2$
leaves the reconstruction unchanged and reduces the penalty. What would training do without the constraint?
\end{exercise}
\begin{exercise}\label{ex:fvu}
In the setting of \cref{prop:shrink}, let $f^\star$ be uniform on $[0,1]$ and $\lambda=0.4$. Compute the mean
squared reconstruction error $\E\,(f^\star-f)^2$ of the optimal $f$.
\end{exercise}
\begin{exercise}\label{ex:ggb}
The ten strongest activations of a latent are all on text about the Golden Gate Bridge. Name two further
observations and one intervention that you would require before calling it a ``Golden Gate Bridge feature''.
\end{exercise}
